\section{Native programs generate the experimental doctrine}\label{sec:source}
The question is not merely how three invariants of one table differ.
An interaction can make a new experiment executable, and a new operation
can turn an old observation into a new experiment by precomposition.
We derive the resulting comparisons from these source operations. There
are two comparison directions and an essential distinction between
strict, lax, and partial transport.

Let $\mathcal C$ be a small category of admitted native programs and
$X:\mathcal C\to\mathsf{Set}$ their deterministic state semantics.
Objects are experimental interfaces. At interface $i$, supply a set
$O_i$ of labelled seed observations, with their actual Boolean response
columns $C_e\subseteq X_i$. A label includes a native observation
witness; a morphism of $\mathcal C$ includes its native program witness.
The functor $X$ need not be faithful: two different native programs may
have the same state transformation.

This is a relative source model, not a theorem that finitely many
observations determine all physical laws or all possible programs.
Actual residual stock, damaged specimens, and installed devices belong
in the source state when they affect subsequent execution. A partial
command can be represented on a \emph{certified domain interface} on
which it is total; this does not turn unsafe disabledness into an
observable outcome. A measurement may have a state-changing native
implementation even though its final Boolean response alone supplies
an incidence column. In particular, intersections of columns below are
semantic conjunctions, not automatically jointly executable tests on
one destructive specimen.

\begin{definition}[Generated experiments]
A generated experiment at $i$ is a labelled pair $(p,e)$ with
$p:i\to j$ and $e\in O_j$, evaluated by
\[
 C_{(p,e)}=X(p)^{-1}(C_e).
\]
Write $G_i$ for the family of its distinct response columns, while
retaining the pairs as native witness annotations. Define
\[
 M_i=\Meet(G_i)
 =\left\{\bigcap\mathcal U:\mathcal U\subseteq G_i\right\},
 \qquad
 \cl_i(A)=\bigcap\{C\in G_i:A\subseteq C\}.
\]
Empty intersections are $X_i$. Define $x\sim_i y$ when all columns of
$G_i$ have the same membership value at $x$ and $y$, and put
$Q_i=X_i/{\sim_i}$.
\end{definition}
The family $G$ is the least family containing the seeds and stable
under preimage along admitted programs. Precomposition acts on
\emph{labels} by $(q,e)\mapsto(qp,e)$, so the semantic construction has
actual native witnesses, not selected endpoint arrows.

\section{Generated state and concept actions}
\begin{theorem}[The source-derived operational doctrine]
For every $p:i\to j$:
\begin{enumerate}
\item $X(p)$ descends uniquely to $\bar p:Q_i\to Q_j$. These maps
make $Q$ a functor, with quotient maps $q_i:X_i\twoheadrightarrow Q_i$
natural in $i$.
\item Preimage restricts to a complete-meet map
$p^*:M_j\to M_i$. Its left adjoint is
\[
 p_!(A)=\cl_j(X(p)[A]).
\]
It preserves all joins and satisfies
$p_!\gamma_i(x)=\gamma_j(X(p)x)$, where
$\gamma_i(x)=\cl_i(\{x\})$.
\item The maps $p_!$ preserve identities and composition. They are the
unique complete-join actions agreeing with native state operations on
all marked object concepts $\gamma_i(x)$.
\end{enumerate}
\end{theorem}
\begin{proof}
A target generated column is $X(r)^{-1}(C_e)$, whose preimage under
$X(p)$ is $X(rp)^{-1}(C_e)\in G_i$. Consequently source-row equality
implies target-row equality. Surjectivity of $q_i$ gives uniqueness and
functoriality of the quotient operations.

Every extent is an intersection of generated columns; preimage preserves
arbitrary intersections, including the empty intersection. It therefore
restricts as claimed. For extents $A\in M_i,B\in M_j$,
\[
 \cl_j(X(p)[A])\subseteq B
 \quad\Longleftrightarrow\quad
 A\subseteq X(p)^{-1}(B),
\]
which proves the adjunction. Since $X(p)^{-1}(\cl_j(X(p)[U]))$ is a
closed set containing $U$, it contains $\cl_i(U)$. Thus
\begin{equation}
 \cl_j(X(p)[\cl_i(U)])=\cl_j(X(p)[U]).\tag{CN}
\end{equation}
This yields the marking equation and the composition law.
Finally $A=\bigvee_{x\in A}\gamma_i(x)$, with joins given by closed
unions. Every complete-join map agreeing on these generators is forced.
\end{proof}

For a one-object category this is a deterministic action of a monoid,
and generated tests are continuations. For a DFA with accepting set $F$,
$G=\{u_w^{-1}(F):w\in\Sigma^*\}$. The quotient is the predictive
row quotient; after retaining only reachable states it is the usual
minimal DFA. The lattice is its FCA/complete-meet completion, not the
state set itself. These fixed-interface facts are established automata,
FCA and Chu-space theory and were already proved in the preceding
\emph{common generation} manuscript.

\section{Growth forces a mixed variance comparison}\label{sec:growth}
Let $J:\mathcal C\hookrightarrow\mathcal D$ extend the admitted
program category, possibly by adding interface objects as well as
programs. Let $Y:\mathcal D\to\mathsf{Set}$ extend the source semantics,
with $YJ=X$. Old seed observations are retained at old interfaces;
new seeds may be supplied at any interface of $\mathcal D$. Decorate
new generated columns, lattices, closures and quotients by $+$.
At every old interface,
\[
 G_i\subseteq G_i^+,\qquad M_i\subseteq M_i^+.
\]
The new experiments can pass through newly introduced interfaces.
They need not be old experiments with renamed labels.

\begin{theorem}[Two directions, and strict backward coherence]
At each old interface the extension induces:
\begin{enumerate}
\item a unique surjection $\rho_i^Q:Q_i^+\twoheadrightarrow Q_i$
with $\rho_i^Qq_i^+=q_i$;
\item a complete-meet inclusion $j_i:M_i\hookrightarrow M_i^+$,
whose left adjoint is the surjective complete-join map
\[
 \rho_i^M:M_i^+\twoheadrightarrow M_i,
 \qquad \rho_i^M(A)=\cl_i(A);
\]
\item strict naturality for every \emph{shared} native program $p:i\to j$:
\begin{align*}
 \rho_j^Q\bar p^+&=\bar p\rho_i^Q,\\
 \rho_j^M p_!^+&=p_!\rho_i^M,\\
 \rho_i^M\gamma_i^+(x)&=\gamma_i(x);
\end{align*}
\item forward lax naturality, in the opposite direction,
\begin{align*}
 p_!^+j_i&\leq j_jp_!,\\
 \gamma_i^+(x)&\subseteq j_i\gamma_i(x).
\end{align*}
All backward comparisons compose strictly under successive extensions.
The forward inequalities are compatible with identities, composition
of shared programs, and successive extensions.
\end{enumerate}
\end{theorem}
\begin{proof}
Equality on more tests implies equality on fewer, proving (1).
For $A\in M_i^+$ and $B\in M_i$,
$\cl_i(A)\subseteq B$ iff $A\subseteq B$, so
$\rho_i^M\dashv j_i$. It is surjective because it fixes $M_i$.
The state square follows by evaluating it on all $x\in X_i$.
On right adjoints there is an equality
$j_i p^*=p^{*+}j_j$, since both sides are the same set-theoretic
inverse image. Taking left adjoints gives the second square. The
marking square follows from
$\cl_i\cl_i^+=\cl_i$.

New closure is smaller: $\cl_j^+(U)\subseteq\cl_j(U)$.
Apply this to $U=X(p)[A]$ to obtain the forward action inequality;
the marking inequality is the singleton case. For two shared programs,
monotonicity gives
\[
 q_!^+p_!^+j_i\leq q_!^+j_jp_!
                  \leq j_kq_!p_!.
\]
These are precisely the unique order comparisons between their
specified endpoints, so they respect the strict program composition
law. For successive growth stages $0\subseteq1\subseteq2$,
$\cl_0\cl_1=\cl_0$ and the inclusions compose literally;
the quotient comparison is unique on source-state representatives.
This proves all stated coherences.
\end{proof}
A newly added program need not act on the old quotient or lattice at all.
The theorem asserts naturality for shared operations, not an impossible
old action for every newly executable operation.

\begin{theorem}[Least simultaneous operational repair]
The generated extent family is the least Moore family containing its
seed columns for which every admitted source program has a
complete-join action agreeing with that program on marked object
concepts. Relative to an old interface, it is the least such extension
of the old extent family and all new seed columns. Independently, its
row equivalence is the largest program-stable equivalence refining the
seed distinctions, hence the coarsest adequate deterministic quotient.
\end{theorem}
\begin{proof}
For any Moore families $N_i$, a source map $f:X_i\to X_j$ has a
marking-compatible complete-join extension iff
\[
 f^{-1}(B)\in N_i\quad\text{for every }B\in N_j.\tag{E}
\]
For necessity, a join extension implies
$f[\cl_{N_i}(U)]\subseteq\cl_{N_j}(f[U])$, because
$\cl_{N_i}(U)$ is the join of the generators of $U$. Set
$U=f^{-1}(B)$: its closure still maps into $B$, so $U$ is closed.
For sufficiency, condition (E) proves (CN) and the adjunction exactly
as above. If every native primitive satisfies (E), induction on
programs forces every generated seed preimage into $N_i$; closure
under intersections forces $M_i\subseteq N_i$. Conversely $M$ itself
satisfies (E).

If an equivalence respects seeds and is stable under every native
program, equivalent states have equal membership in every generated
column, so it refines the constructed row equivalence. The latter is
itself stable by the first theorem. Neither part asserts that the
coarsest row quotient already supplies the least concept repair.
\end{proof}

\section{Three distinct tests for an unchanged characterization}\label{sec:repair}
For clarity take a finite carrier $X$, old columns $F\subseteq\Pow(X)$,
new seed columns $N$, and the \emph{entire resulting} native generator
set $A$, including old operations acting on newly introduced tests.
Old columns can already have been generated by the old operation set,
but that does not make old operations redundant on $N$. Write
\[
 F^+=\Gen_A(F\cup N),\quad M=\Meet(F),\quad M^+=\Meet(F^+).
\]
The object nerve $K_F$ has nonempty faces the nonempty subsets of an
actual column of $F$. Include the empty simplex separately.

\begin{proposition}[Exact repair tests]
\begin{enumerate}
\item $Q_{F^+}\to Q_F$ is bijective iff every new seed is constant on
old row classes and every $a\in A$ respects old row equality.
\item $M=M^+$ iff $N\subseteq M$ and
$a^{-1}(B)\in M$ for all $a\in A,B\in M$. It suffices to check
$B\in F$.
\item $K_F=K_{F^+}$ iff every nonempty new seed and every nonempty
$a^{-1}(C)$, for $a\in A,C\in F$, is contained in an old column.
\end{enumerate}
The second condition implies the first. Neither forces the third.
The third need not imply either of the others.
\end{proposition}
\begin{proof}
For (1), old-class constancy and operation stability propagate through
all words; conversely an unchanged row quotient is stable for all
new operations and observes all new seeds. For (2), inverse image
preserves intersections, so the stated generator check gives closure
of $M$ under all $A$-words. Minimality then gives equality. Necessity
follows from the stability of $M^+$.

For (3), a nonempty column is itself a face because $X$ is finite.
Thus unchanged complexes force each indicated column to be covered.
Conversely let $D(F)$ be the family of subsets covered by old columns,
together with the empty set. The one-step covering condition implies
$a^{-1}(S)\in D(F)$ whenever $S\in D(F)$, by choosing a column
$C\supseteq S$ and using
$a^{-1}(S)\subseteq a^{-1}(C)$. The extra seeds are in $D(F)$;
induction gives $F^+\subseteq D(F)$ and therefore equal nerves.

Every old extent is a union of old row classes, proving the implication.
The CNOT example below separates (1) from (2), and its reset stage
separates (2) from (3). Conversely, an old universal column already
makes the nerve a simplex; adding arbitrary separating columns can
change both rows and the lattice without changing that nerve.
\end{proof}
For infinite carriers, the nerve criterion must say that every finite
subset of each added column is old-covered. A single whole-column
cover need not follow. The exactness assertion here concerns equality
of simplicial complexes, not an equivalence of their homotopy types.

